\documentclass[11pt,a4paper]{article}

\usepackage[margin=1in]{geometry}
\usepackage{amsmath,amssymb}
\usepackage{booktabs}
\usepackage{graphicx}
\usepackage{microtype}
\usepackage[numbers,sort&compress]{natbib}
\usepackage[colorlinks=true,citecolor=blue,urlcolor=blue]{hyperref}

\newcommand{\dd}{\mathrm{d}}
\newcommand{\MSE}{\operatorname{MSE}}
\newcommand{\RMSE}{\operatorname{RMSE}}
\newcommand{\erfc}{\operatorname{erfc}}
\newcommand{\draftnote}[1]{%
  \par\medskip
  \noindent\textbf{Draft status:} #1
  \par\medskip
}

\title{Joint Estimation of Thermal Conductivity and Interfacial
Thermal Resistance in a Layered Slab Using
Physics-Informed Neural Networks}

\author{Virginio Torlao\\
\small Department and institution to be added}
\date{}

\begin{document}

\maketitle

\draftnote{This is a working manuscript.
The introduction and methods describe the implemented benchmark
and the intended comparative study.
Inverse-PINN recovery, repeated training runs, and the complete
noise comparison must be evaluated before publication claims
are added.}

% Write the abstract after the principal results are established.
% Recommended final structure:
% 1. Introduction
% 2. Materials and methods
% 3. Results
% 4. Discussion
% 5. Conclusions
% Include code/data availability and funding declarations.

\section{Introduction}

Heat transport through layered materials depends on both the
thermal properties of the constituent layers and the resistance
to heat flow at their interfaces. A finite interfacial thermal
resistance allows a temperature discontinuity between adjacent
materials even when heat flux is continuous. Consequently,
neglecting imperfect thermal contact can alter the predicted
temperature distribution and the interpretation of transient
thermal measurements. A recent physics-informed neural network
study explicitly addressed the heat equation under imperfect
contact conditions, including the associated temperature jump
and error analysis \citep{oh2025}.

Estimating layer conductivity and interfacial resistance from
temperature observations presents a coupled inverse problem.
Both quantities influence the transient thermal response, and
changes in one parameter may partially compensate for changes
in the other. A close agreement between predicted and observed
temperatures therefore does not, by itself, demonstrate accurate
parameter identification. The reliability of an inverse method
must also be assessed through parameter recovery, sensitivity
to measurement noise, and consistency with the governing
equations and interface conditions.

Physics-informed neural networks (PINNs) combine observations
with residuals of physical equations in a common optimization
problem. They can approximate a state field while simultaneously
estimating unknown coefficients in its governing equations
\citep{raissi2019}. For layered heat conduction, this framework
allows temperature observations to be combined with the heat
equation, flux continuity, and the thermal-resistance jump
condition. However, the representation of the temperature field
requires care: a single globally smooth approximation cannot
directly represent a finite temperature discontinuity at a
material interface.

Domain decomposition provides a way to assign separate neural
networks to different spatial or temporal regions.
Space--time decomposition is an established extension of
physics-informed learning \citep{jagtap2020}.
In the present benchmark, the material interface and the
termination of an imposed heating pulse provide natural
locations for decomposition. The heating pulse also introduces
abrupt changes in the exterior boundary flux, which a smooth
network trained using boundary penalties may approximate
inaccurately.

This study investigates joint identification of the thermal
conductivity of a second material layer and the thermal
resistance of its interface with a first layer. A one-dimensional
synthetic benchmark is used so that the reference parameters,
boundary conditions, and numerical data-generation process
are explicitly controlled. The PINN uses separate temperature
representations for the two layers and for the heating and
redistribution intervals. An analytical boundary lift imposes
the prescribed exterior heat flux, while trainable corrections
represent the finite-domain and interface response.

The intended comparison concerns parameter accuracy,
temperature agreement, physics residuals, sensor placement,
and sensitivity to noise. A conventional numerical
least-squares inversion provides a reference method.
The objective is to establish the conditions under which the
PINN provides reliable identification, rather than to assume
an advantage over conventional inversion.
Because PINNs for imperfect-contact heat conduction already
exist \citep{oh2025}, the contribution must be established
through the completed comparative experiments, not through
the use of a PINN alone.

\section{Materials and methods}

\subsection{Benchmark configuration and assumptions}

The benchmark consists of two homogeneous layers arranged
along the coordinate $x$:
\begin{equation}
\Omega_1=(0,L_1),
\qquad
\Omega_2=(L_1,L_1+L_2).
\end{equation}

Heat transfer is assumed to be one-dimensional.
Each layer has constant thermal conductivity and volumetric
heat capacity. Internal heat generation, phase change,
radiation, lateral heat loss, and temperature dependence of
the material properties are excluded.

The interface has a constant areal thermal resistance $R$.
Its heat capacity is neglected, so it stores no energy.
The exterior right boundary is insulated, and the left
boundary receives a prescribed rectangular heat-flux pulse.

The numerical properties in Table~\ref{tab:benchmark} define
an illustrative benchmark. They do not represent a measured
or specifically identified material pair.

\begin{table}[htbp]
\centering
\caption{Properties and operating conditions of the
synthetic layered-material benchmark.}
\label{tab:benchmark}
\begin{tabular}{lll}
\toprule
Quantity & Symbol & Value \\
\midrule
Layer 1 thickness & $L_1$ & $0.005~\mathrm{m}$ \\
Layer 2 thickness & $L_2$ & $0.005~\mathrm{m}$ \\
Known layer 1 conductivity & $k_1$
  & $10~\mathrm{W\,m^{-1}\,K^{-1}}$ \\
Reference layer 2 conductivity & $k_2^\star$
  & $1~\mathrm{W\,m^{-1}\,K^{-1}}$ \\
Volumetric heat capacity, both layers & $C_v$
  & $2.0\times10^6~\mathrm{J\,m^{-3}\,K^{-1}}$ \\
Reference interface resistance & $R^\star$
  & $10^{-3}~\mathrm{m^2\,K\,W^{-1}}$ \\
Initial temperature & $T_0$ & $293.15~\mathrm{K}$ \\
Imposed pulse flux & $q_0$ & $10^4~\mathrm{W\,m^{-2}}$ \\
Pulse duration & $t_p$ & $10~\mathrm{s}$ \\
Reference simulation duration & $t_{\mathrm{end}}$
  & $200~\mathrm{s}$ \\
Inverse observation interval & & $0.5$--$60~\mathrm{s}$ \\
\bottomrule
\end{tabular}
\end{table}

\subsection{Governing equations}

In layer $i$, the temperature $T_i(x,t)$ satisfies
\begin{equation}
C_v\frac{\partial T_i}{\partial t}
=
k_i\frac{\partial^2T_i}{\partial x^2},
\qquad i=1,2.
\label{eq:heat}
\end{equation}

The heat flux is defined as positive in the direction of
increasing $x$:
\begin{equation}
q_i=-k_i\frac{\partial T_i}{\partial x}.
\label{eq:fourier}
\end{equation}

The initial condition is
\begin{equation}
T_i(x,0)=T_0.
\label{eq:initial}
\end{equation}

At the heated exterior boundary,
\begin{equation}
q_1(0,t)=
\begin{cases}
q_0, & 0<t<t_p,\\
0, & t>t_p.
\end{cases}
\label{eq:leftbc}
\end{equation}

The right exterior boundary is insulated:
\begin{equation}
q_2(L_1+L_2,t)=0.
\label{eq:rightbc}
\end{equation}

At the material interface, conservation of energy requires
\begin{equation}
q_1(L_1^{-},t)=q_2(L_1^{+},t)=q_\Gamma(t).
\label{eq:fluxcontinuity}
\end{equation}

The interface temperature jump is
\begin{equation}
T_1(L_1^{-},t)-T_2(L_1^{+},t)
=
R q_\Gamma(t).
\label{eq:jump}
\end{equation}

Equation~\eqref{eq:jump} is distinct from temporal continuity.
The two material layers may have different temperatures at
their interface, whereas temperature at a fixed location
should remain continuous when the heating pulse terminates.

\subsection{Synthetic reference data}

Reference temperatures are generated using a cell-centered
finite-volume discretization of Eq.~\eqref{eq:heat}.
For adjacent cells $j$ and $j+1$, the thermal conductance
per unit area is
\begin{equation}
G_{j+1/2}
=
\left[
\frac{\Delta x_j}{2k_j}
+
R_{j+1/2}
+
\frac{\Delta x_{j+1}}{2k_{j+1}}
\right]^{-1},
\label{eq:conductance}
\end{equation}
where $R_{j+1/2}=R$ only at the material interface and is
zero at all other internal faces.

The corresponding face flux is
\begin{equation}
q_{j+1/2}=G_{j+1/2}(T_j-T_{j+1}).
\end{equation}

The cell energy balance is
\begin{equation}
C_v\Delta x_j\frac{\dd T_j}{\dd t}
=
q_{j-1/2}-q_{j+1/2}.
\end{equation}

Backward Euler time integration gives
\begin{equation}
C_v\Delta x_j
\frac{T_j^{n+1}-T_j^n}{\Delta t}
=
q_{j-1/2}^{n+1}-q_{j+1/2}^{n+1}.
\end{equation}

The reference configuration uses 80 cells per layer and
$\Delta t=0.025~\mathrm{s}$.
Spatial and temporal refinement checks compare sensor
temperatures across different cell counts and time steps.
These comparisons quantify changes under refinement;
they are not interpreted as exact discretization-error
estimates.

Conservation is checked using the stored energy per unit area,
\begin{equation}
E_{\mathrm{stored}}(t)
=
\sum_j C_v\Delta x_j[T_j(t)-T_0],
\end{equation}
and the imposed energy,
\begin{equation}
E_{\mathrm{input}}(t)=q_0\min(t,t_p).
\end{equation}

For the insulated benchmark, the long-time equilibrium
temperature rise is
\begin{equation}
\Delta T_{\mathrm{eq}}
=
\frac{q_0t_p}{C_v(L_1+L_2)}
=
5~\mathrm{K}.
\end{equation}

A conservative solver can still have discretization error.
Energy conservation and temperature convergence are therefore
evaluated separately.

\subsection{Observation protocol and sensor layouts}

Synthetic observations are extracted at
\begin{equation}
t_n=0.5n~\mathrm{s},
\qquad n=1,\ldots,120.
\end{equation}

Two layouts are considered:
\begin{align}
\mathcal{S}_{\mathrm{same}}
  &=\{2,4\}~\mathrm{mm},\\
\mathcal{S}_{\mathrm{across}}
  &=\{2,7\}~\mathrm{mm}.
\end{align}

The first places both sensors in layer 1.
The second places one sensor in each layer.
Both layouts contain the same number of observations.

The current PINN pilot uses $\mathcal{S}_{\mathrm{across}}$.
The complete layout comparison requires running the PINN
with both layouts under the same training protocol.

Noisy observations are defined by
\begin{equation}
y_{s,n}
=
T(x_s,t_n;k_2^\star,R^\star)
+
\varepsilon_{s,n},
\qquad
\varepsilon_{s,n}\sim
\mathcal{N}(0,\sigma_T^2).
\label{eq:noise}
\end{equation}

The conventional-inversion pilot uses
$\sigma_T=0.01~\mathrm{K}$ and 30 noise realizations.
The shared 2~mm sensor receives the same noise realization
in both layouts. The remaining sensor noises are independently
drawn.

For a fair method comparison, the PINN and conventional
inversion must use identical noisy observations.
The completed conventional noise pilot does not constitute
a completed PINN noise study.

\subsection{Dimensionless formulation}

For the equal-thickness layers, define local coordinates
\begin{equation}
z_1=\frac{x}{L_1},
\qquad
z_2=\frac{x-L_1}{L_2},
\qquad
\tau=\frac{t}{t_p}.
\end{equation}

The dimensionless temperature rise is
\begin{equation}
u_i=\frac{T_i-T_0}{T_s},
\qquad T_s=5~\mathrm{K}.
\end{equation}

Let
\begin{equation}
\kappa_i=\frac{k_i}{k_{\mathrm{ref}}},
\qquad
k_{\mathrm{ref}}=1~\mathrm{W\,m^{-1}\,K^{-1}}.
\end{equation}

The dimensionless heat equation becomes
\begin{equation}
\frac{\partial u_i}{\partial\tau}
-
a_i\frac{\partial^2u_i}{\partial z_i^2}
=0,
\qquad
a_i=\frac{k_it_p}{C_vL_i^2}.
\end{equation}

For the benchmark,
\begin{equation}
a_1=2,
\qquad
a_2=0.2\kappa_2.
\end{equation}

The normalized rightward flux is
\begin{equation}
\widehat q_i
=
\frac{q_i}{q_0}
=
-\frac{k_iT_s}{L_iq_0}
\frac{\partial u_i}{\partial z_i}.
\end{equation}

Thus,
\begin{equation}
\widehat q_1=-u_{1,z},
\qquad
\widehat q_2=-0.1\kappa_2u_{2,z}.
\end{equation}

The interface conditions become
\begin{align}
\widehat q_1(1,\tau)
  &=\widehat q_2(0,\tau),\\
u_1(1,\tau)-u_2(0,\tau)
  &=\gamma\widehat q_1(1,\tau),
\qquad
\gamma=\frac{Rq_0}{T_s}.
\end{align}

At the reference resistance, $\gamma=2$.

\subsection{Spatial and temporal network decomposition}

Each material layer has two neural networks: one for
$0\leq\tau<1$ and one for $1\leq\tau\leq6$.
The implementation therefore contains four networks.

Each network has two inputs, three hidden layers containing
32 neurons each, and one scalar output.
Hyperbolic tangent activation functions are used.
Weights are initialized using Xavier normal initialization,
with zero initial biases.

The time inputs are
\begin{equation}
\eta_H=2\tau-1,
\qquad
\eta_C=2\frac{\tau-1}{5}-1,
\end{equation}
for heating and redistribution, respectively.

The spatial inputs are
\begin{equation}
\xi_1=2z_1^2-1,
\qquad
\xi_2=2(1-z_2)^2-1.
\label{eq:spatialfeatures}
\end{equation}

These transformations make the spatial derivative of the
learned correction zero at the left exterior of layer 1
and the right exterior of layer 2.
They do not prescribe the temperature at either exterior.

The heating corrections are
\begin{equation}
v_{i,H}(z,\tau)
=
g(\tau)\mathcal{N}_{i,H}(\xi_i,\eta_H),
\qquad
g(\tau)=\frac{\tau}{\tau+0.1}.
\end{equation}

The redistribution corrections are
\begin{equation}
v_{i,C}(z,\tau)
=
\mathcal{N}_{i,C}(\xi_i,\eta_C).
\end{equation}

The heating gate enforces a zero initial correction.
The connection between the two time intervals is imposed
through a temperature-continuity loss, rather than an exact
constraint.

\subsection{Analytical lift for the imposed heat-flux pulse}

The left exterior flux is represented using a known solution
of the dimensionless half-line heat equation.
For $s>0$, define
\begin{equation}
F(z,s)
=
2\sqrt{\frac{a_1s}{\pi}}
\exp\left(-\frac{z^2}{4a_1s}\right)
-
z\,\erfc\left(\frac{z}{2\sqrt{a_1s}}\right),
\label{eq:stepresponse}
\end{equation}
and set $F(z,s)=0$ for $s\leq0$.

At regular points,
\begin{equation}
F_s-a_1F_{zz}=0,
\end{equation}
and
\begin{equation}
F_z(0,s)=-1,
\qquad s>0.
\end{equation}

The pulse lift is
\begin{equation}
B(z,\tau)=F(z,\tau)-F(z,\tau-1).
\label{eq:pulselift}
\end{equation}

The predicted temperatures are
\begin{equation}
u_1(z,\tau)
=
B(z,\tau)+
\begin{cases}
v_{1,H}(z,\tau), & \tau<1,\\
v_{1,C}(z,\tau), & \tau\geq1,
\end{cases}
\label{eq:layer1representation}
\end{equation}
and
\begin{equation}
u_2(z,\tau)
=
\begin{cases}
v_{2,H}(z,\tau), & \tau<1,\\
v_{2,C}(z,\tau), & \tau\geq1.
\end{cases}
\label{eq:layer2representation}
\end{equation}

Equations~\eqref{eq:spatialfeatures}--%
\eqref{eq:layer2representation} enforce the imposed exterior
flux and right insulation away from the pulse corner times.
The initial temperature is also satisfied.

The half-line solution is used only as a boundary lift.
It is not the solution of the finite two-layer problem.
The learned corrections must still satisfy the finite-domain
heat equations and material-interface conditions.

The lift depends on the known conductivity $k_1$.
It does not use the unknown $k_2$ or $R$.

Classical boundary derivatives are not assessed at the
incompatible corner $(z_1,\tau)=(0,0)$ or at the exact flux
switch $(0,1)$. Temperature predictions at the switch remain
part of the observation protocol.

\subsection{Trainable material properties}

The unknown properties are represented by logarithmic
parameters:
\begin{equation}
k_2=k_{\mathrm{ref}}\exp(\alpha),
\qquad
R=R_{\mathrm{ref}}\exp(\beta),
\label{eq:logparameters}
\end{equation}
where
$R_{\mathrm{ref}}=1~\mathrm{m^2\,K\,W^{-1}}$.

This parameterization guarantees positivity.
It does not impose finite upper or lower bounds.

The initial inverse-PINN guesses are
\begin{equation}
k_{2,\mathrm{init}}
=
0.5~\mathrm{W\,m^{-1}\,K^{-1}},
\qquad
R_{\mathrm{init}}
=
0.002~\mathrm{m^2\,K\,W^{-1}}.
\end{equation}

Both properties are optimized jointly with the network
weights and biases.

During fixed-parameter diagnostics, the properties are held
at their reference values.
Those diagnostics still use sensor data in the loss.
They are therefore data-assisted field-fitting tests,
not observation-free forward solutions.

\subsection{Loss function}

For sampled values $\{r_j\}_{j=1}^{N}$, define
\begin{equation}
\MSE(r)=\frac{1}{N}\sum_{j=1}^{N}r_j^2.
\end{equation}

The PDE residuals are
\begin{equation}
f_1=u_{1,\tau}-2u_{1,zz},
\qquad
f_2=u_{2,\tau}-0.2\kappa_2u_{2,zz}.
\end{equation}

The PDE loss is
\begin{equation}
\mathcal{L}_{\mathrm{PDE}}
=
\MSE(f_1)+\MSE(f_2).
\end{equation}

At the material interface, define
\begin{align}
r_q&=\widehat q_1(1,\tau)-\widehat q_2(0,\tau),\\
r_R&=u_1(1,\tau)-u_2(0,\tau)
-\gamma\widehat q_1(1,\tau).
\end{align}

The interface losses are
\begin{equation}
\mathcal{L}_q=\MSE(r_q),
\qquad
\mathcal{L}_R=\MSE(r_R).
\end{equation}

For dimensionless observations
$\widetilde y_{s,n}=(y_{s,n}-T_0)/T_s$, the data loss is
\begin{equation}
\mathcal{L}_{\mathrm{data}}
=
\sum_{s=1}^{2}
\MSE\left[
u_{i(s)}(z_s,\tau_n)-\widetilde y_{s,n}
\right],
\end{equation}
where $i(s)$ identifies the layer containing sensor $s$.

The temporal connection loss is
\begin{equation}
\mathcal{L}_{\mathrm{time}}
=
\sum_{i=1}^{2}
\MSE\left[
v_{i,C}(z,1)-v_{i,H}(z,1)
\right].
\end{equation}

The analytical lift is shared by the two time branches and
therefore cancels in this difference.

The implemented total loss is
\begin{equation}
\mathcal{L}
=
\mathcal{L}_{\mathrm{PDE}}
+\mathcal{L}_q
+\mathcal{L}_R
+10\mathcal{L}_{\mathrm{data}}
+10\mathcal{L}_{\mathrm{time}}.
\label{eq:totalloss}
\end{equation}

Exterior boundary residuals are retained as diagnostics but
excluded from Eq.~\eqref{eq:totalloss}, because the exterior
flux conditions are imposed by construction.

Spatial and temporal derivatives are evaluated using
automatic differentiation.
All calculations use double precision.

\subsection{Collocation sampling and optimization}

Half of the interior collocation points are sampled uniformly
in the heating interval $0<\tau<1$, and half in the
redistribution interval $1<\tau<6$.
Spatial coordinates are sampled uniformly in $0<z<1$.
A small numerical exclusion avoids the exact pulse corner
times.

This sampling gives equal numbers of points to intervals
of unequal duration.
The resulting loss is not a uniformly time-weighted integral
over the complete observation interval.

The Adam stage uses 512 interior points per layer and
256 shared time samples for evaluating the material-interface
conditions. Collocation points are refreshed every 100 steps.
Temporal continuity is evaluated at 129 uniformly spaced
spatial points in each layer.

Training starts with 5000 Adam steps at learning rate
$10^{-3}$.
It is followed by L-BFGS with a maximum of 500 iterations,
history size 50, and a strong-Wolfe line search.

During L-BFGS, the collocation set is fixed and contains
1024 interior points per layer and 512 interface time samples.
The number of closure evaluations may exceed the number of
optimizer iterations.

The pilot uses random seed 42 and an NVIDIA GeForce
RTX 4050 Laptop GPU.
Final reporting should include the installed software
versions and results from repeated independent training seeds.

\subsection{Conventional numerical inversion}

A conventional least-squares inverse solver provides a
comparison method.
Its forward model uses the finite-volume spatial
discretization with 60 cells per layer.

The semidiscrete system can be written as
\begin{equation}
\frac{\dd\boldsymbol{\vartheta}}{\dd t}
=
-\mathbf{A}\boldsymbol{\vartheta}
+\mathbf{b}q_{\mathrm{in}}(t),
\end{equation}
where $\boldsymbol{\vartheta}$ contains cell temperature
rises.

Modal decomposition of $\mathbf{A}$ gives the transient
response for the prescribed rectangular pulse without
backward-Euler time stepping.
Sensor temperatures are obtained by interpolation from the
cell values.

The estimated parameters minimize
\begin{equation}
\Phi(k_2,R)
=
\sum_{s,n}
\left[
T_{\mathrm{num}}(x_s,t_n;k_2,R)-y_{s,n}
\right]^2.
\end{equation}

Optimization uses logarithmic parameters with bounds
\begin{equation}
0.1\leq k_2\leq10
\quad
[\mathrm{W\,m^{-1}\,K^{-1}}],
\end{equation}
and
\begin{equation}
10^{-5}\leq R\leq10^{-2}
\quad
[\mathrm{m^2\,K\,W^{-1}}].
\end{equation}

The clean-data pilot uses three starting pairs:
\begin{equation}
(k_2,R)
\in
\left\{
(0.5,0.002),\,
(2,0.0002),\,
(5,0.005)
\right\},
\end{equation}
with the units given above.

The conventional and PINN methods currently differ in their
parameter bounds.
This difference must be disclosed in comparisons of failure
rates and robustness.

The reference generator and conventional inverse solver also
use different spatial resolutions and time integration.
Their discretization mismatch can introduce estimation bias,
including for clean synthetic observations.

\subsection{Accuracy, sensitivity, and physics diagnostics}

Parameter recovery is measured using relative errors:
\begin{equation}
e_{k_2}
=
100\frac{\widehat k_2-k_2^\star}{k_2^\star},
\qquad
e_R
=
100\frac{\widehat R-R^\star}{R^\star}.
\end{equation}

Temperature agreement is quantified by
\begin{equation}
\RMSE_T
=
\sqrt{
\frac{1}{2N_t}
\sum_{s=1}^{2}\sum_{n=1}^{N_t}
[\widehat T(x_s,t_n)-y_{s,n}]^2
}.
\end{equation}

Because the current observations are used during training,
this quantity measures agreement with fitted data.
It is not a held-out prediction error.

Physics diagnostics use a fresh sample containing 2048
interior points per layer and 1024 interface time samples.
PDE, interface-flux, interface-temperature-jump, and data
residuals are reported separately.

The temporal connection is additionally evaluated on a
513-point spatial grid:
\begin{equation}
\Delta T_i^{\mathrm{time}}(z)
=
T_s[v_{i,C}(z,1)-v_{i,H}(z,1)].
\end{equation}

Its RMS and maximum absolute values are reported for each
layer.

Boundary-flux diagnostics are evaluated separately during
startup, the heating plateau, the pulse switch, and subsequent
redistribution.
For the direct-flux representation, small exterior boundary
errors are an architectural consequence and are not treated
as independent evidence of successful field reconstruction.

Local parameter sensitivity is assessed using a Jacobian
with respect to the logarithmic parameters:
\begin{equation}
J_{m\ell}
=
\frac{\partial T_m}{\partial p_\ell},
\qquad
\boldsymbol{p}
=
\left(
\ln\frac{k_2}{k_{\mathrm{ref}}},
\ln\frac{R}{R_{\mathrm{ref}}}
\right).
\end{equation}

Singular values, the Jacobian condition number, and the cosine
between its two sensitivity columns characterize local
parameter coupling.
These quantities provide local diagnostics and do not
establish global identifiability.

For repeated noise realizations, parameter bias, standard
deviation, and RMS estimation error are reported separately.
A reduction in total estimation error is not interpreted as
a reduction in noise sensitivity unless the spread also
decreases.

\subsection{Exploratory development and remaining validation}

The development sequence examined a single time network,
additional sampling near pulse transitions, temporal
decomposition, an increased temporal-continuity weight,
and the analytical exterior-flux lift.

These were exploratory changes made after inspecting pilot
diagnostics. They are not presented as a preregistered
comparison or as independent evidence of a generally optimal
architecture.

Total loss values are not directly compared across
configurations that change the sampling distribution,
loss weights, or included terms.
Common temperature errors and separately defined physics
diagnostics provide the appropriate comparisons.

Before final conclusions are drawn, the study requires
repeated inverse-PINN runs, matched noisy-data comparisons,
both sensor layouts, and additional checks against reference
temperatures away from the fitted sensors.
Integrated energy balance and material-interface predictions
should also be evaluated.
A successful single clean-data inverse run is treated as a
pilot result rather than a complete robustness demonstration.

\section{Results}

% Suggested subsection order:
% 3.1 Reference-solver convergence and energy conservation
% 3.2 Fixed-parameter PINN diagnostics and exploratory changes
% 3.3 Clean-data recovery of k_2 and R
% 3.4 Dependence on initialization and training seed
% 3.5 Sensor-layout comparison and local sensitivity
% 3.6 Noise robustness and comparison with numerical inversion
%
% Keep parameter-recovery results separate from temperature-fit results.
% Do not populate inverse-PINN claims before the runs are completed.

\section{Discussion}

% Discuss:
% - Why direct flux enforcement improved the observed pilot fit.
% - Remaining interface and PDE errors.
% - Parameter coupling and layout-dependent sensitivity.
% - Numerical-data mismatch and estimation bias.
% - Computational cost and training variability.
% - What the method adds relative to the numerical inverse solver.
% - Limits of a synthetic, constant-property, one-dimensional benchmark.
% - Relationship to existing imperfect-contact PINN literature.

\section{Conclusions}

% Write after the comparative study is complete.
% State only conclusions supported by the completed experiments.

\section*{Data and code availability}

% Add the actual repository URL, release tag, and archived data location
% once the reproducibility package has been prepared.

\begin{thebibliography}{99}

\bibitem{raissi2019}
M.~Raissi, P.~Perdikaris, and G.~E.~Karniadakis,
``Physics-informed neural networks: A deep learning framework
for solving forward and inverse problems involving nonlinear
partial differential equations,''
\emph{Journal of Computational Physics},
vol.~378, pp.~686--707, 2019.
\href{https://doi.org/10.1016/j.jcp.2018.10.045}
{doi:10.1016/j.jcp.2018.10.045}.

\bibitem{jagtap2020}
A.~D.~Jagtap and G.~E.~Karniadakis,
``Extended physics-informed neural networks (XPINNs):
A generalized space-time domain decomposition based deep
learning framework for nonlinear partial differential
equations,''
\emph{Communications in Computational Physics},
vol.~28, no.~5, pp.~2002--2041, 2020.
\href{https://doi.org/10.4208/cicp.OA-2020-0164}
{doi:10.4208/cicp.OA-2020-0164}.

\bibitem{oh2025}
H.~Oh and G.~Jo,
``Physics-informed neural network for the heat equation
under imperfect contact conditions and its error analysis,''
\emph{AIMS Mathematics},
vol.~10, no.~4, pp.~7920--7940, 2025.
\href{https://doi.org/10.3934/math.2025364}
{doi:10.3934/math.2025364}.

\end{thebibliography}

\end{document}